\documentclass[11pt,a4paper]{article}
\usepackage{microtype}
\usepackage[margin=2.6cm]{geometry}
\usepackage{amsmath,amsthm,thmtools,thm-restate}
\usepackage{fontspec}
\usepackage{unicode-math}
\directlua{luaotfload.add_fallback("cfb", {"NotoSansMath:mode=harf;", "DejaVuSans:mode=harf;"})}
\setmainfont{Libertinus Serif}[RawFeature={fallback=cfb}]
\setmathfont{Latin Modern Math}
\setmonofont{Latin Modern Mono}[Scale=0.85,RawFeature={fallback=cfb}]
\AtBeginDocument{\let\setminus\smallsetminus}
\usepackage{enumitem}
\usepackage{longtable,booktabs,array,calc}
\usepackage{tcolorbox}
\usepackage{fancyhdr}
\usepackage[colorlinks=true,linkcolor=blue!50!black,citecolor=blue!50!black,urlcolor=blue!50!black]{hyperref}
\usepackage[capitalize,nameinlink]{cleveref}

\theoremstyle{plain}
\newtheorem{theorem}{Theorem}[section]
\newtheorem{lemma}[theorem]{Lemma}
\newtheorem{corollary}[theorem]{Corollary}
\newtheorem{proposition}[theorem]{Proposition}
\newtheorem{observation}[theorem]{Observation}
\theoremstyle{definition}
\newtheorem{definition}[theorem]{Definition}
\newtheorem{question}[theorem]{Question}
\newtheorem{conjecture}[theorem]{Conjecture}
\theoremstyle{remark}
\newtheorem{remark}[theorem]{Remark}
\newtheorem*{proofidea}{Idea of proof}
\newcommand{\fullproofref}[1]{\,\hyperref[#1]{\textit{[full proof]}}}
\providecommand{\tightlist}{\setlength{\itemsep}{0pt}\setlength{\parskip}{0pt}}

\newcommand{\E}{\mathbb E}
\newcommand{\chif}{\chi_{\mathrm{f}}}
\newcommand{\Hrand}{H_{\mathrm{rand}}}
\newcommand{\ind}{\mathbf 1}
\DeclareMathOperator{\Bin}{Bin}
\newcommand{\repair}[1]{\begin{tcolorbox}[colback=blue!4,colframe=blue!40!black,boxrule=0.4pt,arc=2pt,left=4pt,right=4pt,top=2pt,bottom=2pt]\small\textbf{Repair.} #1\end{tcolorbox}}

\pagestyle{fancy}\fancyhf{}
\fancyhead[L]{\small\itshape Candidate result a\_generalization\_of\_vizings\_theorem --- machine-generated proof, awaiting human review}
\fancyhead[R]{\small\thepage}
\renewcommand{\headrulewidth}{0.2pt}

\title{A counterexample to Rosenfeld's hypergraph generalization\\ of Vizing's theorem}
\author{\href{https://github.com/graph-theory-AI}{github.com/graph-theory-AI}\\[2pt] \normalsize Machine-generated note}
\date{17 September 2026}

\begin{document}
\maketitle

\begin{abstract}
Rosenfeld conjectured (Open Problem Garden, 2007) that every simple $d$-uniform hypergraph in which each $(d-1)$-set lies in at most $r$ edges has an $(r+d-1)$-edge-colouring in which any two edges sharing $d-1$ vertices receive distinct colours; for $d=2$ this is Vizing's theorem. We show that the conjecture fails for $d=100$, $r=200$: a probabilistic construction, whose main ingredient is an entropy bound on the number of large matchings in a regular linear hypergraph, yields a simple $100$-uniform hypergraph with maximum $99$-codegree at most $200$ that needs at least $451$ colours, against the conjectured $299$. The same method gives, for every $d\ge15$, hypergraphs with maximum codegree $2d$ needing $(1-o(1))\,d^2/(5\log d)$ colours, so the conjectured bound $r+d-1$ is off by an unbounded factor.
The counterexample is non-constructive and astronomically large: the proof produces it only for $n\ge10^{456}$ vertices per class, i.e.\ about $10^{458}$ vertices and $10^{45\,000}$ edges. It cannot be instantiated or spot-checked, and the disproof rests entirely on the correctness of the counting lemma, every ingredient of which was verified separately by the referee on small instances.

\medskip\noindent\small
This note was produced by AI within the \href{https://github.com/graph-theory-AI/Graph-Theory-LLM-Proofs}{Graph Theory AI} project:
the argument was found by a language model and audited by an adversarial LLM
referee, and it has not been checked by a human mathematician.
\Cref{sec:provenance} records the full provenance.
\end{abstract}

\section{Introduction}

A hypergraph $H$ is \emph{$d$-uniform} if every edge has exactly $d$ vertices, and \emph{simple} if it has no repeated edges (so $E(H)\subseteq\binom{V(H)}{d}$). A \emph{facet} of an edge $e$ is one of its $d$ subsets of size $d-1$; the \emph{codegree} of a $(d-1)$-set is the number of edges containing it. A \emph{facet colouring} of $H$ is a map $E(H)\to C$ such that two distinct edges sharing a facet, i.e.\ sharing $d-1$ vertices, receive different colours, and $\chif(H)$ denotes the least number of colours in a facet colouring. All logarithms are natural.

The problem posed by Rosenfeld on the Open Problem Garden \cite{OPG} reads, verbatim:

\begin{conjecture}[{Rosenfeld \cite{OPG}}]\label{conj:rosenfeld}
Let $H$ be a simple $d$-uniform hypergraph, and assume that every set of $d-1$ points is contained in at most $r$ edges. Then there exists an $r+d-1$-edge-coloring so that any two edges which share $d-1$ vertices have distinct colors.
\end{conjecture}

In the notation above the conjecture states $\chif(H)\le r+d-1$ whenever every $(d-1)$-set has codegree at most $r$. For $d=2$ the hypothesis says $\Delta(H)\le r$ and the conclusion is $\chi'(H)\le r+1$, which is Vizing's theorem \cite{Vizing64}. The OPG entry remarks that Kempe-chain arguments, the main tool for Vizing's theorem, do not seem to work for $d\ge3$; apart from a 2009 comment asking for a reference, the entry carries no further information, and the referee's literature search (\cref{app:referee}) found no paper attacking the conjecture.

\paragraph{The facet hypergraph.} It is convenient to translate the problem. Given $H$, let $F=F(H)$ be the hypergraph whose vertices are the $(d-1)$-subsets of $V(H)$ that are facets of some edge, with one edge $\{$the $d$ facets of $e\}$ for each $e\in E(H)$. Then $F$ is $d$-uniform, its maximum degree is the maximum codegree of $H$, and $F$ is \emph{linear} (two distinct edges share at most one vertex), because two distinct $d$-sets cannot have two common $(d-1)$-subsets. A facet colouring of $H$ is exactly a proper edge colouring of $F$, so \cref{conj:rosenfeld} asserts $\chi'(F)\le\Delta(F)+d-1$ for every linear $d$-uniform $F$ that is \emph{realisable} as a facet hypergraph. Realisability is essential: for general linear hypergraphs the bound is false by a wide margin (a projective plane of order $q$ is a linear $(q+1)$-uniform hypergraph with $\Delta=q+1$ and $\chi'=q^2+q+1$), but a projective plane is not a facet hypergraph, since $d$-sets pairwise meeting in $d-1$ points form a star or lie inside a $(d+1)$-set. The construction below respects this constraint by building an actual $H$.

\paragraph{Related work.} The conjecture is sharp: the referee computed exact facet chromatic numbers of complete $d$-uniform hypergraphs $\binom{[d+k]}{d}$ (here $r=k+1$ and $r+d-1=d+k$) and found the bound attained with equality at $(d,d+k)=(2,3),(2,5),(3,4),(3,5),(3,6),(4,5),(5,6),(5,7),(6,7),(7,8)$, and never exceeded; $2357$ random instances with $d\in\{3,4,5\}$ on at most $d+4$ vertices gave no violation either. Any counterexample therefore has to be large and sparse, which is what the probabilistic construction produces.

Several neighbouring lines of work concern edge colourings of linear hypergraphs in other regimes and are not in conflict with our result. Pippenger and Spencer \cite{PS89} and Kahn \cite{Kahn96} proved $\chi'(F)=(1+o(1))\Delta(F)$ for linear $k$-uniform $F$ with $k$ fixed and $\Delta\to\infty$; our hypergraphs have $k=d$ and $\Delta\le2d$, so $k$ is of the order of $\Delta$ and these theorems do not apply (and projective planes show they cannot be extended to that regime). Alon and Kim \cite{AK97} study $t$-simple uniform hypergraphs, again for fixed uniformity and large degree. The Berge--F\"uredi conjecture \cite{Furedi86,Berge89}, $\chi'(F)\le\Delta([F]_2)+1$ for linear $F$, bounds the chromatic index by the maximum degree of the \emph{shadow graph} rather than of $F$, and has seen recent progress \cite{BFH24a,BFH24b,MA25}; Kang, Kelly, K\"uhn, Methuku and Osthus \cite{KKKMO21} solved a problem of Erd\H{o}s on the (ordinary) chromatic index of hypergraphs with bounded codegree. None of these addresses colourings in which only edges sharing $d-1$ vertices must differ. To the best of the referee's search, the counterexample below is new; see \cref{rem:novelty}.

\paragraph{The result.} Fix $d=100$ and let $n$ be a multiple of $5$. Take disjoint classes $V_1,\dots,V_{100}$ of size $n$, let $H_0$ be the complete $100$-partite $100$-uniform hypergraph (all transversal $100$-sets), retain each edge independently with probability $100/n$, and delete every retained edge containing a facet of codegree above $200$. The resulting hypergraph $H$ is simple, $100$-uniform, and has maximum $99$-codegree at most $200$. A first-moment count of large matchings in the facet hypergraph shows that, with overwhelming probability, no colour class can contain $n^{99}/5$ edges, while $H$ keeps more than $90\,n^{99}$ edges; hence $\chif(H)\ge451$.

\begin{restatable}[Main theorem]{theorem}{thmmain}\label{thm:main}
Let $n\ge10^{456}$ be a multiple of $5$. There exists a simple, $100$-partite, $100$-uniform hypergraph $H$ with $100n$ vertices such that
\begin{enumerate}[nosep]
\item every $99$-set of vertices is contained in at most $200$ edges, and
\item every facet colouring of $H$ uses at least $451$ colours.
\end{enumerate}
\end{restatable}

\begin{corollary}\label{cor:main}
\Cref{conj:rosenfeld} is false for $d=100$ and $r=200$: the conjectured bound is $r+d-1=299<451\le\chif(H)$.
\end{corollary}
\begin{proof}
Immediate from \cref{thm:main}.
\end{proof}

\begin{proofidea}
Let $T=n^{99}$. The facet hypergraph $F$ of $H_0$ has $N=100T$ vertices, is $100$-uniform, $n$-regular and linear, and a set of edges of $H_0$ is a legal colour class iff the corresponding edges of $F$ form a matching. \Cref{lem:count} bounds the number $a_m(F)$ of $m$-matchings of a regular linear hypergraph from above; for $m=T/5$ and sampling probability $p=100/n$ it gives $\E[\#\text{$m$-matchings in the sampled $F$}]\le e^{-T/5}$ (\cref{prop:firstmoment}), so with probability at least $1-e^{-T/5}$ every colour class of the sampled hypergraph $\Hrand$, and of every subhypergraph, has fewer than $T/5$ edges. Pruning the edges through overloaded facets removes only an expected $100(e/4)^{100}$ fraction of the edges (Markov's inequality applied to $2^{X}$, $X\sim\Bin(n-1,100/n)$), so some outcome has both properties and more than $90T$ edges; $450$ colours would cover fewer than $450\cdot T/5=90T$ edges (\cref{prop:prune}).\fullproofref{pf:main}
\end{proofidea}

The exponent in the first-moment count improves as $d$ grows, and the same argument shows that the conjectured bound is wrong by an unbounded factor.

\begin{restatable}[Large uniformity]{theorem}{thmlarge}\label{thm:large}
For every integer $d\ge15$ there exists a simple $d$-uniform hypergraph $H_d$ in which every $(d-1)$-set lies in at most $2d$ edges and
\[
\chif(H_d)\ \ge\ \bigl(1-2d\,(e/4)^{d}\bigr)\frac{d^2}{\lceil5\log d\rceil}.
\]
Consequently $\chif(H_d)/(r+d-1)\ge(1-o(1))\,d/(15\log d)\to\infty$ as $d\to\infty$, where $r=2d$.
\end{restatable}

\begin{proofidea}
Repeat the construction with $d$ classes of size $n$, sampling probability $d/n$, pruning threshold $2d$ and matching size $xT$ with $x=\lceil5\log d\rceil/d$, $T=n^{d-1}$. For fixed $d\ge15$ the first-moment exponent is negative, so letting $n\to\infty$ (through multiples of $d$) gives an outcome with at least $dT(1-2d(e/4)^d)$ edges and no colour class of size $xT$. The limits are taken in the order ``fix $d$, then $n\to\infty$''; no uniformity in $d$ is claimed or needed.\fullproofref{pf:large}
\end{proofidea}

\section{Counting matchings in regular linear hypergraphs}\label{sec:count}

The counting lemma is proved by the entropy method, in the manner of Radhakrishnan's proof of Br\'egman's theorem \cite{Bregman73,Radhakrishnan97} and of the Kahn--Lov\'asz theorem \cite{AF08,CR11}: a uniformly random perfect matching is revealed vertex by vertex in a random order, and the entropy of each reveal is bounded by the logarithm of the total weight of the still-available edges. The hypergraph version needs two modifications: edge weights, and a small allowance $\varepsilon$ for edges that meet a matching edge in two or more vertices.

\begin{restatable}[Weighted perfect matchings]{lemma}{lemPM}\label{lem:pm}
Let $s\ge2$ and let $K$ be an $s$-uniform hypergraph on $L$ vertices with positive edge weights $w$, such that every vertex has weighted degree $\sum_{e\ni v}w(e)\le\Delta$. Put $Z=\sum_{M}\prod_{e\in M}w(e)$, the sum over perfect matchings $M$ of $K$. Call an edge $e$ \emph{bad relative to $M$} if $|e\cap f|\ge2$ for some $f\in M$; note that every edge of $M$ is bad relative to $M$. Suppose that for every perfect matching $M$ and every vertex $v$,
\begin{equation}\label{eq:badhyp}
\sum_{e\ni v,\ e\text{ bad relative to }M}w(e)\ \le\ \varepsilon\Delta .
\end{equation}
Then
\begin{equation}\label{eq:pm}
\log Z\ \le\ \frac Ls\Bigl[\log\Delta+\int_0^1\log\bigl(u^{s-1}+\varepsilon\bigr)\,du\Bigr].
\end{equation}
\end{restatable}

\begin{proofidea}
Let $M$ be a perfect matching drawn with probability $\prod_{e\in M}w(e)/Z$, so that $\log Z=H(M)+\E\sum_{e\in M}\log w(e)$. Give the vertices independent uniform labels in $[0,1]$ and reveal $M(v)$, the matching edge at $v$, in decreasing order of labels. By the chain rule and Gibbs' inequality, $\log Z\le\sum_v\E[\ind\{v\text{ has the largest label in }M(v)\}\log S_v]$, where $S_v$ is the total weight of the edges at $v$ disjoint from all previously revealed matching edges. Given the label $t$ of $v$, the event has probability $t^{s-1}$, and a non-bad edge at $v$ is still available with probability $t^{s(s-1)}$ because its other $s-1$ vertices lie in $s-1$ distinct matching edges other than $M(v)$. Jensen's inequality and the substitution $u=t^s$ give \eqref{eq:pm}.\fullproofref{pf:pm}
\end{proofidea}

As $\varepsilon\to0$ the integral in \eqref{eq:pm} tends to $\int_0^1(s-1)\log u\,du=-(s-1)$, which is the form in which the original writeup used the lemma. The rate matters, and we make it explicit.

\begin{restatable}[Quantitative form of the integral]{lemma}{lemInt}\label{lem:int}
For every integer $s\ge3$ and every $\varepsilon\in(0,1)$,
\[
\int_0^1\log\bigl(u^{s-1}+\varepsilon\bigr)\,du\ \le\ -(s-1)+(s+1)\,\varepsilon^{1/(s-1)} .
\]
\end{restatable}

\repair{The referee's only substantive objection was that the writeup asserts $\int_0^1\log(u^{s-1}+\varepsilon)\,du=-(s-1)+o(1)$ without a rate, whereas the true correction is $\approx(s-1)\varepsilon^{1/(s-1)}$, a $99$th root for $s=100$, and it is multiplied by $L/s\approx80\,n^{99}$ against a margin of only $0.2\,n^{99}$ in \cref{prop:firstmoment}. \Cref{lem:int} is the missing quantification; its proof splits the integral at $u_0=\varepsilon^{1/(s-1)}$ and uses $\log(1+y)\le y$ above $u_0$. The constant $s+1$ is not optimal ($s-1$ is the asymptotically correct one) but suffices.}

\begin{proofidea}
Below $u_0=\varepsilon^{1/(s-1)}$ the integrand is at most $\log(2\varepsilon)$; above $u_0$ it is at most $(s-1)\log u+\varepsilon u^{-(s-1)}$, and $\int_{u_0}^1\varepsilon u^{-(s-1)}du\le u_0/(s-2)$. The $\log\varepsilon$ terms cancel exactly.\fullproofref{pf:int}
\end{proofidea}

\begin{restatable}[Matchings in a regular linear hypergraph]{lemma}{lemCount}\label{lem:count}
Let $s\ge3$ and $x\in(0,1)$. Let $F$ be an $s$-uniform, $D$-regular, linear hypergraph on $N$ vertices such that $m=xN/s$ is an integer, and put $u=N-sm=(1-x)N$ and $b=(s-1)u$. Assume $b\ge4$ and
\[
\varepsilon:=\frac xD+\frac{s^3}{b}<1 .
\]
Then the number $a_m(F)$ of matchings of size $m$ in $F$ satisfies
\begin{equation}\label{eq:count}
\log a_m(F)\ \le\ m\bigl(\log(D/x)-(s-1)\bigr)-N(1-x)\log(1-x)+(m+u)(s+1)\,\varepsilon^{1/(s-1)} .
\end{equation}
\end{restatable}

\repair{The writeup states \eqref{eq:count} with an unquantified $o(N)$ in place of the last term, arising from three sources: the $o(1)$ of the integral, an $O(N^{-1})$ in $\varepsilon$ from bad auxiliary edges, and Stirling's formula. Here all three are explicit: the integral by \cref{lem:int}, the bad-edge fraction by a direct count giving $s^3/b$, and Stirling is replaced by the one-sided bounds $\log b!\ge b\log b-b$ and $\binom{b}{s-1}\le b^{s-1}/(s-1)!$, which are all that the argument needs and carry no error term. The referee verified the auxiliary construction, the weighted regularity and the identity \eqref{eq:Z} exactly on five small instances (\cref{app:referee}).}

\begin{proofidea}
Adjoin a set $B$ of $b$ new vertices and, for every old vertex $a$ and every $C\in\binom{B}{s-1}$, an auxiliary edge $\{a\}\cup C$ of weight $w=\frac{(1-x)D/x}{\binom{b}{s-1}}$; the original edges keep weight $1$. The resulting $K$ is weighted-regular of degree $\Delta=D/x$, and its perfect matchings correspond to pairs ($m$-matching of $F$, partition of $B$ into $(s-1)$-blocks assigned to the $u$ uncovered old vertices), so
\begin{equation}\label{eq:Z}
Z=a_m(F)\cdot\frac{b!}{((s-1)!)^u}\cdot w^u .
\end{equation}
Linearity of $F$ and the shape of the auxiliary edges make hypothesis \eqref{eq:badhyp} hold with the stated $\varepsilon$. \Cref{lem:pm,lem:int} bound $\log Z$ from above, the one-sided Stirling bounds bound the completion factor in \eqref{eq:Z} from below, and subtraction gives \eqref{eq:count}.\fullproofref{pf:count}
\end{proofidea}

\section{The random hypergraph}\label{sec:random}

Throughout this section $d=s=100$, $n$ is a multiple of $5$, $T=n^{99}$, $V_1,\dots,V_{100}$ are disjoint classes of size $n$, and $H_0$ is the hypergraph of all transversal $100$-sets, so $|E(H_0)|=n^{100}$. Its facet hypergraph $F$ has as vertices the $100\,n^{99}=100T$ transversal $99$-sets; each lies in exactly $n$ edges of $H_0$, so $F$ is $100$-uniform, $n$-regular and linear on $N=100T$ vertices. Retain each edge of $H_0$ independently with probability $p=100/n$; call the result $\Hrand$. A set of edges of $\Hrand$ is a legal colour class iff the corresponding edges of $F$ form a matching.

\begin{restatable}[No large colour class]{proposition}{propFM}\label{prop:firstmoment}
Let $m=T/5$ and let $A$ be the event that $\Hrand$ contains no $m$ edges that are pairwise facet-disjoint. If $n\ge10^{456}$ then $\Pr(A^{c})\le e^{-T/5}$.
\end{restatable}

\repair{The writeup proves $\log\E[\#\text{$m$-matchings}]\le-\tfrac25T+o(N)$ and concludes $\le-\tfrac15T$ ``for all sufficiently large $n$''. With \cref{lem:count} the error term is $(m+u)(s+1)\varepsilon^{1/99}=80.2\,T\cdot101\cdot\varepsilon^{1/99}$ with $\varepsilon\le1/(4n)$, and it is at most the available margin $0.2\,T$ exactly when $n\ge10^{455.6}$. The threshold $10^{456}$ in \cref{thm:main} is where this happens; the referee's independent estimate was $n\gtrsim10^{454.6}$ (with the slightly smaller constant $s-1$). Using the exact value $-0.7056$ of the exponent instead of $-0.4$ would lower the threshold to about $10^{416}$, still far beyond anything that can be written down.}

\begin{proofidea}
The expected number of $m$-matchings in the sampled facet hypergraph is $p^m a_m(F)$. Apply \cref{lem:count} with $x=1/5$, $D=n$: the main term is $T\bigl[x\log(100/x)+x-100f(x)\bigr]$ with $f(x)=x+(1-x)\log(1-x)\ge x^2/2$, hence at most $T[(\log500+1)/5-2]<-\tfrac25T$; the error term is at most $8100.2\,T\,(4n)^{-1/99}\le0.198\,T$ for $n\ge10^{456}$. Markov's inequality finishes.\fullproofref{pf:firstmoment}
\end{proofidea}

\begin{restatable}[Pruning]{proposition}{propPrune}\label{prop:prune}
Delete from $\Hrand$ every edge containing a $99$-set of codegree greater than $200$, and let $H$ be the result and $Z=|E(H)|$. Then every $99$-set of vertices lies in at most $200$ edges of $H$, $Z\le200T$ always, and $\E Z\ge100T\bigl(1-100(e/4)^{100}\bigr)>99T$. Consequently, if $n\ge10^{456}$ is a multiple of $5$, some outcome satisfies both $A$ and $Z>90T$.
\end{restatable}

\begin{proofidea}
Conditional on an edge being present, the number of other sampled edges through one of its facets is $\Bin(n-1,100/n)$, and Markov's inequality on $2^X$ gives $\Pr(X\ge200)\le2^{-200}(1+100/n)^{n-1}\le(e/4)^{100}=:q\approx1.7\cdot10^{-17}$; a union bound over the $100$ facets gives $\E Z\ge100T(1-100q)$. The deterministic bound $Z\le200T$ follows by summing facet degrees. Then $\E[Z\ind_A]\ge99T-200T\,e^{-T/5}>90T$.\fullproofref{pf:prune}
\end{proofidea}

\section{Remarks}\label{sec:remarks}

\begin{remark}[The counterexample cannot be exhibited]\label{rem:size}
\Cref{thm:main} is an existence statement about an object with $100n\ge10^{458}$ vertices and about $90\,n^{99}\approx10^{45\,150}$ edges. Nothing of this size can be constructed, stored or spot-checked, and the disproof therefore rests entirely on the correctness of \cref{lem:pm,lem:int,lem:count}. The referee, unable to test the counterexample itself, tested every ingredient instead: inequality \eqref{eq:pm} by exhaustive enumeration on $95$ weighted hypergraphs with $s\in\{2,3,4,5\}$ and at most $12$ vertices (no violation; typical near-tight cases $K_{6,6}$ with $\log Z=6.579$ against the bound $7.621$, and $\mathrm{AG}(2,3)$ with $1.386$ against $2.150$); the auxiliary construction of \cref{lem:count} exactly on $C_4,C_6,K_4,K_{3,3}$ and $\mathrm{AG}(2,3)$; and the constant $-(s-1)$ in \eqref{eq:count} independently, as the first moment of the number of $m$-matchings in a random $D$-regular $s$-uniform hypergraph, which agrees with the right-hand side of \eqref{eq:count} to within $10^{-4}$ relative for $s\in\{2,3,5,100\}$. The error term in \eqref{eq:count} is genuinely needed: on $K_N$ with $x=1$ the bound without it fails by an additive $0.847$ at every $N$ tested up to $10^8$. This is why the referee's confidence is medium rather than high, and why a second expert reading of \cref{lem:pm} is warranted before any claim of priority.
\end{remark}

\begin{remark}[Sharpness of the conjecture and size of any counterexample]\label{rem:sharp}
The bound $r+d-1$ is attained with equality on complete $d$-uniform hypergraphs for the ten pairs $(d,|V|)$ listed in the introduction, and the referee's $2357$ random instances with $d\in\{3,4,5\}$ and $|V|\le d+4$ produced no violation, with minimum slack $0$. So the conjecture is sharp and non-trivial, no small counterexample exists in that range, and a counterexample had to be a large sparse object of the kind produced here. Nothing in this note concerns $d=3$ or other small uniformities, or small $r$; \cref{conj:rosenfeld} remains open there.
\end{remark}

\begin{remark}[Consistency with known theorems; the numbers]\label{rem:consistency}
The mechanism is a first-moment count showing that a random sparse $d$-partite hypergraph has no matching covering a fifth of the facets, so that $\chif\ge|E(H)|/(T/5)\approx5d$. It does not conflict with the Pippenger--Spencer and Kahn theorems \cite{PS89,Kahn96}, which require the uniformity $k$ to be fixed while $\Delta\to\infty$; here $k=d=100$ and $\Delta\le200$. The referee recomputed every constant: the first-moment exponent at $x=1/5$ is exactly $-0.7056$ (the writeup's weakened bound is $-0.557<-0.4$); $q=(e/4)^{100}=1.67\cdot10^{-17}$; and $450(m-1)=90T-450<90T$. The method is not pushed to its limit: the first-moment exponent vanishes at $x^{\ast}=0.1435$, so the same argument with $|E(H)|>99T$ gives $\chif\ge690$, and a random-greedy heuristic (which covers a fraction $0.0887$ of the facets) suggests the true value is nearer $1100$.
\end{remark}

\begin{remark}[Two points a reader may stumble on]\label{rem:bad}
(i)~In \cref{lem:pm} an edge of $M$ is bad relative to $M$, since $|e\cap e|=s\ge2$; this is deliberate and necessary, because $S_v\ge w(M(v))$ forces $\varepsilon\ge w(M(v))/\Delta$, and a reader who excludes $e=f$ obtains a false inequality (a single edge has $Z=1$ against a ``bound'' of $e^{-(s-1)}$). (ii)~The error term in \cref{lem:count} depends on $F$ only through $(N,D,s,x)$; the writeup's claim that it is uniform over $F$ is thereby justified, although the application does not need it, as $F$ is fixed for each $n$.
\end{remark}

\begin{remark}[Novelty]\label{rem:novelty}
The referee found no counterexample to \cref{conj:rosenfeld} and no paper attacking it, in the OPG entry (fetched live; the only post-2007 content is a 2009 request for a reference) or in the neighbouring literature cited in the introduction. Novelty was checked only against what could be found online, and the counting lemma is a routine variant of a known entropy argument; what is new, if anything, is its application to Rosenfeld's problem.
\end{remark}

\section{Provenance}\label{sec:provenance}

The mathematics in this note was produced without human intervention, in three stages.
(1)~The construction and proof were found by the language model \texttt{gpt-6-astra} (reasoning effort max, single autonomous pass) on 2026-09-16, as part of a sweep over the open problems of the Open Problem Garden (catalog id \texttt{a\_generalization\_of\_vizings\_theorem}, leg \texttt{attacks\_opg}; original writeup in \texttt{attacks\_opg/a\_generalization\_of\_vizings\_theorem/output.md}). The model claimed the verdict ``disproved'' with high confidence.
(2)~An independent adversarial referee agent (\texttt{claude-opus-5}) reviewed the writeup on 2026-09-17. It re-derived every step, brute-forced inequality \eqref{eq:pm} on $95$ instances and the auxiliary construction on five, confirmed the constant $-(s-1)$ independently, recomputed all arithmetic, tested the conjecture itself on complete and on $2357$ random small hypergraphs, fetched the OPG entry and searched for prior art. Its verdict was \textsc{minor\_gaps} with confidence \emph{medium}: no step was wrong, but the $o(N)$ term of the counting lemma was never quantified, and its quantification (which the referee carried out) shows that the counterexample exists only for $n\gtrsim10^{455}$. Its report is reproduced verbatim in \cref{app:referee}.
(3)~On 2026-09-17 the writeup was rewritten into the present note by \texttt{claude-fable-5-1}. The text was reorganised, with full proofs in \cref{app:proofs}, and the following changes were made, each marked where it occurs:
\begin{itemize}[nosep]
\item \Cref{lem:int} (explicit rate $(s+1)\varepsilon^{1/(s-1)}$ for the integral) was added; it replaces the writeup's unquantified limit (4). The constant $s+1$, the bad-edge bound $s^3/b$ and the threshold $10^{456}$ are the rewriter's own explicit choices, made along the lines of the referee's computation; the referee did not see this exact form.
\item \Cref{lem:count} is stated with an explicit error term: $\varepsilon=x/D+s^3/b$ (the writeup has $\Delta^{-1}+O_{s,x}(N^{-1})$, with a sketched union bound that is here carried out), and the Stirling step (9) of the writeup is replaced by the one-sided bounds $\log b!\ge b\log b-b$, $\binom b{s-1}\le b^{s-1}/(s-1)!$, which remove its $o(N)$.
\item \Cref{prop:firstmoment} and \cref{thm:main} state the explicit threshold $n\ge10^{456}$ where the writeup says ``sufficiently large $n$''; the abstract and \cref{rem:size} state the size of the resulting object.
\item The two-sentence entropy step (6) of the writeup is written out in the proof of \cref{lem:pm}, following the referee's expansion (chain rule conditional on the labels, ``$v$ uncovered'' $=$ ``$v$ has the largest label in $M(v)$'', each edge charged once, Gibbs' inequality).
\item The writeup's \S4 is stated as \cref{thm:large} with the order of limits made explicit (fix $d\ge15$, then $n\to\infty$) and the constant $1-2d(e/4)^d$ in place of $1-o_d(1)$; the writeup mixed $o_d(1)$ and $o_n(1)$.
\item The observation that every edge of $M$ is bad relative to $M$ is stated in \cref{lem:pm} and explained in \cref{rem:bad}, as the referee recommended.
\item The introduction, the citations \cite{Vizing64,PS89,Kahn96,AK97,Furedi86,Berge89,BFH24a,BFH24b,MA25,KKKMO21,Bregman73,Radhakrishnan97,AF08,CR11}, the facet-hypergraph reformulation with the projective-plane example, and the contents of \cref{sec:remarks} are taken from the referee report; the writeup cites nothing and calls its counting lemma new. The entropy method is attributed to Radhakrishnan and to the proofs of the Kahn--Lov\'asz theorem.
\end{itemize}
No other mathematical content was changed.
\textbf{No human mathematician has checked this argument.} We circulate this proof for human review, not as a claim to a new resolution.

\appendix

\section{Full proofs}\label{app:proofs}

\phantomsection\label{pf:pm}
\lemPM*
\begin{proof}
If $Z=0$ there is nothing to prove, so assume $K$ has a perfect matching. Let $M$ be a random perfect matching with $\Pr(M)=\prod_{e\in M}w(e)/Z$. Taking expectations in $-\log\Pr(M)=\log Z-\sum_{e\in M}\log w(e)$ gives
\begin{equation}\label{eq:HZ}
H(M)+\E\sum_{e\in M}\log w(e)=\log Z .
\end{equation}

\emph{The reveal process.} Independently of $M$, give each vertex $y$ a label $\ell_y$, the $\ell_y$ being i.i.d.\ uniform on $[0,1]$; almost surely they are distinct. For a vertex $v$ write $M(v)$ for the edge of $M$ containing $v$, and say that $v$ is \emph{first in $M(v)$} if $\ell_v=\max_{y\in M(v)}\ell_y$. Given $M$ and the labels, let $R_v=\{M(y):\ell_y>\ell_v\}$ be the set of matching edges containing a vertex of larger label than $v$, and let
\[
S_v=\sum\{\,w(e):\ e\ni v,\ e\cap f=\emptyset\text{ for all }f\in R_v\,\}
\]
be the total weight of the edges at $v$ that avoid all edges of $R_v$. If $v$ is first in $M(v)$ then $M(v)\notin R_v$ and $M(v)$ is disjoint from every edge of $R_v$, so $S_v\ge w(M(v))>0$.

\emph{Chain rule.} Since the labels $\ell$ are independent of $M$, $H(M)=H(M\mid\ell)$. Given $\ell$, list the vertices $v_1,v_2,\dots,v_L$ in decreasing order of label. The matching $M$ is determined by the sequence $(M(v_1),\dots,M(v_L))$, so by the chain rule
\[
H(M\mid\ell)=\sum_{i=1}^L H\bigl(M(v_i)\ \big|\ \ell,\ M(v_1),\dots,M(v_{i-1})\bigr).
\]
Fix $i$ and a value of the conditioning, i.e.\ the labels and the edges $M(v_1),\dots,M(v_{i-1})$; these edges are exactly the elements of $R_{v_i}$. If $v_i$ lies in one of them, then $v_i$ is not first in $M(v_i)$ and $M(v_i)$ is determined, so the $i$th term vanishes. Otherwise $v_i$ is first in $M(v_i)$, and since $M$ is a matching, the conditional distribution $P$ of $M(v_i)$ is supported on the set $\mathcal A$ of edges $e\ni v_i$ disjoint from all edges of $R_{v_i}$, whose total weight is $S_{v_i}$. Gibbs' inequality, in the form
\[
H(P)+\sum_{e\in\mathcal A}P(e)\log w(e)\ \le\ \log\sum_{e\in\mathcal A}w(e)
\]
(the difference of the two sides is $-D(P\,\|\,Q)\le0$ for $Q(e)=w(e)/\sum_{\mathcal A}w$), gives
\[
H\bigl(M(v_i)\mid\text{conditioning}\bigr)+\E\bigl[\log w(M(v_i))\mid\text{conditioning}\bigr]\ \le\ \log S_{v_i}
\]
whenever $v_i$ is first in $M(v_i)$. Averaging over the conditioning and summing over $i$,
\[
H(M)+\E\sum_{v}\ind\{v\text{ first in }M(v)\}\log w(M(v))\ \le\ \sum_v\E\bigl[\ind\{v\text{ first in }M(v)\}\log S_v\bigr].
\]
Every edge $f\in M$ has exactly one vertex that is first in $f$, so the sum on the left equals $\sum_{f\in M}\log w(f)$, and by \eqref{eq:HZ}
\begin{equation}\label{eq:six}
\log Z\ \le\ \sum_v\E\bigl[\ind\{v\text{ first in }M(v)\}\log S_v\bigr].
\end{equation}

\emph{Bounding one term.} Fix a vertex $v$ and a perfect matching $M$, and condition on $M$ and on $\ell_v=t$. The other $s-1$ vertices of $M(v)$ have independent labels, so $\Pr(v\text{ first in }M(v)\mid M,t)=t^{s-1}$. Condition further on this event, which involves only the labels of the vertices of $M(v)$. Let $e\ni v$ be an edge that is not bad relative to $M$. Then $|e\cap f|\le1$ for all $f\in M$; in particular $e\cap M(v)=\{v\}$, and the other $s-1$ vertices $y$ of $e$ lie in $s-1$ pairwise distinct matching edges $M(y)$, all different from $M(v)$. The edge $e$ contributes to $S_v$ iff none of these $M(y)$ belongs to $R_v$, i.e.\ iff all $s(s-1)$ vertices of $\bigcup_{y\in e\setminus\{v\}}M(y)$ have labels below $t$. These vertices are disjoint from $M(v)$, so their labels are independent of the conditioning, and the probability is $t^{s(s-1)}$. The bad edges at $v$ have total weight at most $\varepsilon\Delta$ by \eqref{eq:badhyp}. Hence
\[
\E\bigl[S_v\ \big|\ M,\ \ell_v=t,\ v\text{ first in }M(v)\bigr]\ \le\ \Delta t^{s(s-1)}+\varepsilon\Delta .
\]
Since $S_v>0$ on the event, Jensen's inequality for the concave function $\log$ gives
\[
\E\bigl[\ind\{v\text{ first in }M(v)\}\log S_v\ \big|\ M\bigr]
=\int_0^1 t^{s-1}\,\E\bigl[\log S_v\mid M,t,\text{first}\bigr]\,dt
\ \le\ \int_0^1 t^{s-1}\log\bigl(\Delta(t^{s(s-1)}+\varepsilon)\bigr)\,dt .
\]
This bound does not depend on $M$ or $v$. Summing over the $L$ vertices in \eqref{eq:six} and substituting $u=t^s$ (so $du=s\,t^{s-1}dt$ and $t^{s(s-1)}=u^{s-1}$),
\[
\log Z\ \le\ L\int_0^1 t^{s-1}\log\bigl(\Delta(t^{s(s-1)}+\varepsilon)\bigr)\,dt
=\frac Ls\int_0^1\log\bigl(\Delta(u^{s-1}+\varepsilon)\bigr)\,du
=\frac Ls\Bigl[\log\Delta+\int_0^1\log(u^{s-1}+\varepsilon)\,du\Bigr]. \qedhere
\]
\end{proof}

\phantomsection\label{pf:int}
\lemInt*
\begin{proof}
Put $u_0=\varepsilon^{1/(s-1)}\in(0,1)$, so that $u_0^{s-1}=\varepsilon$ and $(s-1)\log u_0=\log\varepsilon$.

For $0\le u\le u_0$ we have $u^{s-1}+\varepsilon\le2\varepsilon$, hence
\[
\int_0^{u_0}\log(u^{s-1}+\varepsilon)\,du\ \le\ u_0\log(2\varepsilon)=u_0\log2+u_0\log\varepsilon .
\]
For $u_0\le u\le1$, write $\log(u^{s-1}+\varepsilon)=(s-1)\log u+\log(1+\varepsilon u^{-(s-1)})\le(s-1)\log u+\varepsilon u^{-(s-1)}$. Now
\begin{align*}
\int_{u_0}^1(s-1)\log u\,du
&=(s-1)\bigl[u\log u-u\bigr]_{u_0}^1\\
&=-(s-1)+(s-1)u_0-(s-1)u_0\log u_0\\
&=-(s-1)+(s-1)u_0-u_0\log\varepsilon,
\end{align*}
and, as $s\ge3$,
\begin{align*}
\int_{u_0}^1\varepsilon u^{-(s-1)}\,du
&=\frac{\varepsilon\bigl(u_0^{-(s-2)}-1\bigr)}{s-2}
 \le \frac{\varepsilon u_0^{-(s-2)}}{s-2}\\
&=\frac{u_0^{\,s-1-(s-2)}}{s-2}=\frac{u_0}{s-2}.
\end{align*}
Adding the three estimates, the terms $\pm u_0\log\varepsilon$ cancel and
\begin{align*}
\int_0^1\log(u^{s-1}+\varepsilon)\,du
&\le -(s-1)+u_0\left[(s-1)+\log2+\frac1{s-2}\right]\\
&\le -(s-1)+(s+1)u_0,
\end{align*}
since $\log2+1/(s-2)\le\log2+1<2$ for $s\ge3$.
\end{proof}

\phantomsection\label{pf:count}
\lemCount*
\begin{proof}
Recall $u=N-sm=(1-x)N$ and $b=(s-1)u$.

\emph{The auxiliary hypergraph.} Adjoin to $V(F)$ a set $B$ of $b$ new vertices, and let $K$ be the weighted $s$-uniform hypergraph on $V(F)\cup B$ ($L=N+b=s(m+u)$ vertices) with the edges of $F$, each of weight $1$, and for every old vertex $a\in V(F)$ and every $C\in\binom{B}{s-1}$ the auxiliary edge $\{a\}\cup C$ of weight
\[
w=\frac{(1-x)D/x}{\binom{b}{s-1}} .
\]
An old vertex has weighted degree $D+\binom b{s-1}w=D+(1-x)D/x=D/x$. A new vertex lies in $N\binom{b-1}{s-2}$ auxiliary edges, and $\binom{b-1}{s-2}/\binom{b}{s-1}=(s-1)/b$, so its weighted degree is $N\frac{s-1}{b}\cdot\frac{(1-x)D}{x}=\frac{(s-1)(1-x)N D}{(s-1)(1-x)N\,x}=D/x$. Thus $K$ is weighted-regular with $\Delta=D/x$.

\emph{Perfect matchings of $K$.} Let $M$ be a perfect matching of $K$. Each auxiliary edge contains exactly $s-1$ new vertices and each original edge none, so the auxiliary edges of $M$ partition $B$ and there are exactly $b/(s-1)=u$ of them; the remaining $(N-u)/s=m$ edges of $M$ are original and form an $m$-matching of $F$, whose $u$ uncovered old vertices are exactly the old vertices of the auxiliary edges. Conversely, every $m$-matching of $F$ extends to a perfect matching of $K$ by choosing an ordered partition of $B$ into $u$ blocks of size $s-1$, one for each uncovered old vertex, in $b!/((s-1)!)^u$ ways. Hence
\[
Z=a_m(F)\cdot\frac{b!}{((s-1)!)^u}\cdot w^u ,
\]
which is \eqref{eq:Z}.

\emph{Hypothesis \eqref{eq:badhyp}.} Fix a perfect matching $M$ and a vertex $v$. An edge $e$ is bad relative to $M$ iff two distinct vertices of $e$ lie in a common edge of $M$.

Original edges at $v$ (so $v$ is old): an auxiliary edge of $M$ contains only one old vertex and so meets an original edge in at most one vertex; an original edge of $M$ meets a different original edge in at most one vertex by linearity of $F$. So the only original edge at $v$ that can be bad is $M(v)$ itself, of weight $1$. The bad original weight at $v$ is at most $1$.

Auxiliary edges at $v$: they all have the same weight $w$, so their total bad weight is (weighted auxiliary degree of $v$)$\times$(fraction of them that are bad) $\le\Delta\cdot\beta$, where $\beta$ is the probability that a uniformly random auxiliary edge at $v$ is bad. We bound $\beta$ by a union bound over the $\binom s2$ pairs of positions in the edge. Note that an edge of $M$ containing a new vertex is auxiliary and so contains exactly $s-1$ new vertices and one old vertex.

\emph{Case $v$ old.} The random edge is $\{v\}\cup C$ with $C$ uniform in $\binom B{s-1}$; list $C$ in uniformly random order as $c_1,\dots,c_{s-1}$, so that each $c_j$ is uniform on $B$ and each pair $(c_i,c_j)$, $i\ne j$, is a uniformly random ordered pair of distinct elements of $B$. For a pair $(v,c_j)$: $\Pr(c_j\in M(v))=|M(v)\cap B|/b\le(s-1)/b$. For a pair $(c_i,c_j)$: $\Pr(c_j\in M(c_i))=|M(c_i)\cap B\setminus\{c_i\}|/(b-1)=(s-2)/(b-1)$. Hence
\[
\beta\ \le\ \frac{(s-1)^2}{b}+\binom{s-1}{2}\frac{s-2}{b-1}.
\]
\emph{Case $v$ new.} The random edge is $\{a\}\cup C$ with $a$ uniform on $V(F)$, $v\in C$, and $C\setminus\{v\}$ uniform in $\binom{B\setminus\{v\}}{s-2}$, independently of $a$. For a pair $(a,c)$ with $c\in C$: given $C$, $a$ is uniform on $V(F)$ and $M(c)$ has exactly one old vertex, so $\Pr(a\in M(c))=1/N$; there are $s-1$ such pairs. For a pair $(v,c)$ with $c\in C\setminus\{v\}$: $c$ is uniform on $B\setminus\{v\}$, so $\Pr(c\in M(v))=(s-2)/(b-1)$; there are $s-2$ such pairs. For a pair $(c,c')$ of distinct elements of $C\setminus\{v\}$: it is a uniformly random ordered pair of distinct elements of $B\setminus\{v\}$, so $\Pr(c'\in M(c))\le(s-2)/(b-2)$; there are $\binom{s-2}2$ such pairs. Since $N=b/((s-1)(1-x))\ge b/(s-1)$, we get $(s-1)/N\le(s-1)^2/b$, and
\[
\beta\ \le\ \frac{(s-1)^2}{b}+\frac{(s-2)^2}{b-1}+\binom{s-2}{2}\frac{s-2}{b-2}\ \le\ \frac{(s-1)^2}{b}+\binom{s-1}{2}\frac{s-2}{b-2},
\]
because $(s-2)^2+\binom{s-2}2(s-2)=(s-2)^2\bigl(1+\tfrac{s-3}{2}\bigr)=\tfrac{(s-2)^2(s-1)}2=\binom{s-1}2(s-2)$.

In both cases, using $b\ge4$ (so $b\le2(b-2)$),
\[
\beta\ \le\ \frac{(s-1)^2+\tfrac12(s-1)(s-2)^2}{b-2}\ \le\ \frac{\tfrac12(s-1)s^2}{b-2}\ \le\ \frac{s^3}{2(b-2)}\ \le\ \frac{s^3}{b},
\]
where the second inequality is $(s-1)+\tfrac12(s-2)^2\le\tfrac12s^2$, i.e.\ $2s-2+s^2-4s+4\le s^2$, i.e.\ $2\le2s$. Altogether the bad weight at $v$ is at most $1+\Delta s^3/b=\Delta\bigl(x/D+s^3/b\bigr)=\varepsilon\Delta$, which is \eqref{eq:badhyp}.

\emph{Upper bound on $\log Z$.} By \cref{lem:pm} and then \cref{lem:int} (applicable since $s\ge3$ and $0<\varepsilon<1$), with $L/s=m+u$ and $\Delta=D/x$,
\begin{equation}\label{eq:Zup}
\log Z\ \le\ (m+u)\Bigl[\log\frac Dx-(s-1)+(s+1)\varepsilon^{1/(s-1)}\Bigr].
\end{equation}

\emph{Lower bound on the completion factor.} Since $\log b!=\sum_{k=1}^b\log k\ge\int_0^b\log t\,dt=b\log b-b$ and $\binom b{s-1}\le b^{s-1}/(s-1)!$,
\begin{align*}
\log\Bigl(\frac{b!}{((s-1)!)^u}w^u\Bigr)
&=\log b!-u\log(s-1)!+u\log\frac{(1-x)D}{x}-u\log\binom b{s-1}\\
&\ge\bigl(b\log b-b\bigr)-u\log(s-1)!+u\log\frac{(1-x)D}{x}-u\bigl((s-1)\log b-\log(s-1)!\bigr)\\
&=-b+u\log\frac{(1-x)D}{x}
= u\Bigl[\log\frac{(1-x)D}{x}-(s-1)\Bigr],
\end{align*}
using $u(s-1)\log b=b\log b$ and $b=(s-1)u$.

\emph{Conclusion.} Subtracting the last bound from \eqref{eq:Zup} in \eqref{eq:Z},
\begin{align*}
\log a_m(F)&\le(m+u)\Bigl[\log\frac Dx-(s-1)\Bigr]-u\Bigl[\log\frac{(1-x)D}{x}-(s-1)\Bigr]+(m+u)(s+1)\varepsilon^{1/(s-1)}\\
&=m\Bigl[\log\frac Dx-(s-1)\Bigr]+u\log\frac1{1-x}+(m+u)(s+1)\varepsilon^{1/(s-1)},
\end{align*}
and $u\log\frac1{1-x}=-N(1-x)\log(1-x)$, which is \eqref{eq:count}.
\end{proof}

\phantomsection\label{pf:firstmoment}
\propFM*
\begin{proof}
Colour classes of $\Hrand$ correspond to matchings of the random edge-subhypergraph $F'$ of $F$ obtained by retaining each edge independently with probability $p=100/n$. A fixed $m$-matching of $F$ survives in $F'$ with probability $p^m$, so the expected number of $m$-matchings of $F'$ is $p^m a_m(F)$, and by Markov's inequality $\Pr(A^c)\le p^m a_m(F)$.

Apply \cref{lem:count} with $s=100$, $x=1/5$, $D=n$, $N=100T$, $m=xN/s=T/5$ (an integer as $5\mid n$), $u=(1-x)N=80T$, $b=99u=7920T$, $m+u=80.2T$. Here $\varepsilon=x/D+s^3/b=\frac1{5n}+\frac{10^6}{7920\,n^{99}}$. For $n\ge2$ the second term is at most $\frac1{20n}$ (as $n^{98}\ge2^{98}>2600$), so $\varepsilon\le\frac1{4n}<1$; also $b\ge4$. Then
\begin{align*}
\log\bigl(p^m a_m(F)\bigr)
&\le m\log p+m\bigl(\log(n/x)-99\bigr)
       -N(1-x)\log(1-x)\\
&\qquad +(m+u)\cdot101\cdot\varepsilon^{1/99}\\
&=\frac N{100}\left[x\log\frac{100}{x}+x-100f(x)\right]
       +80.2\,T\cdot101\cdot\varepsilon^{1/99},
\end{align*}
where $f(x)=x+(1-x)\log(1-x)$. We used $m\log(pn/x)=m\log(100/x)$ and
\[
m(-99)-N(1-x)\log(1-x)
=\frac N{100}\bigl[x-100x-100(1-x)\log(1-x)\bigr]
=\frac N{100}\bigl[x-100f(x)\bigr].
\]

\emph{Main term.} $f(0)=f'(0)=0$ and $f''(x)=1/(1-x)\ge1$, so $f(x)\ge x^2/2$ and $100f(1/5)\ge2$. With $\log500<7$,
\[
x\log\frac{100}x+x-100f(x)\ \le\ \frac{\log500+1}{5}-2\ <\ \frac85-2=-\frac25,
\]
so the main term is at most $-\frac25\cdot\frac N{100}=-\frac25T$. (The exact value of the bracket is $-0.7056$.)

\emph{Error term.} For $n\ge10^{456}$, $\varepsilon^{1/99}\le(4n)^{-1/99}\le\bigl(4\cdot10^{456}\bigr)^{-1/99}=10^{-(456+\log_{10}4)/99}<10^{-4.612}<2.45\cdot10^{-5}$, so the error term is at most $80.2\cdot101\cdot2.45\cdot10^{-5}\,T<0.199\,T\le\frac15T$.

Hence $\log\E[\#\text{$m$-matchings of }F']\le-\frac25T+\frac15T=-\frac T5$, and $\Pr(A^c)\le e^{-T/5}$.
\end{proof}

\phantomsection\label{pf:prune}
\propPrune*
\begin{proof}
Every $99$-set of $V(H)$ that is not a transversal of $99$ classes lies in no edge of $H_0$; every transversal $99$-set has codegree at most $200$ in $H$ by construction. So the codegree condition holds for all $99$-sets.

\emph{Expected size.} Fix an edge $e$ of $H_0$ and one of its facets $f$. The facet $f$ lies in exactly $n$ edges of $H_0$, namely $e$ and $n-1$ others, each retained independently with probability $100/n$. Conditional on $e$ being retained, the number $X$ of other retained edges through $f$ is $\Bin(n-1,100/n)$, and $f$ has codegree greater than $200$ iff $X\ge200$. By Markov's inequality applied to $2^X$,
\[
\Pr(X\ge200)\le2^{-200}\,\E 2^X=2^{-200}\Bigl(1+\frac{100}n\Bigr)^{n-1}\le\frac{e^{100}}{2^{200}}=\Bigl(\frac e4\Bigr)^{100}=:q ,
\]
and $q<10^{-16}$. A union bound over the $100$ facets of $e$ shows that $e$ survives in $H$ with probability at least $p(1-100q)$, so
\[
\E Z\ \ge\ n^{100}\cdot\frac{100}n\,(1-100q)=100T(1-100q)>99T .
\]

\emph{Deterministic bound.} Summing codegrees over the $100T$ transversal $99$-sets counts every edge of $H$ exactly $100$ times, so $100Z\le200\cdot100T$ and $Z\le200T$.

\emph{A good outcome.} Let $n\ge10^{456}$ be a multiple of $5$. By \cref{prop:firstmoment} and the two bounds just proved,
\[
\E[Z\ind_A]\ \ge\ \E Z-200T\Pr(A^c)\ >\ 99T-200T\,e^{-T/5}\ >\ 90T,
\]
the last step because $T=n^{99}\ge16>5\log(200/9)\approx15.5$. If every outcome in $A$ had $Z\le90T$ then $\E[Z\ind_A]\le90T$; so some outcome satisfies $A$ and $Z>90T$.
\end{proof}

\phantomsection\label{pf:main}
\thmmain*
\begin{proof}
Let $n\ge10^{456}$ be a multiple of $5$ and fix an outcome satisfying $A$ and $|E(H)|>90T$, which exists by \cref{prop:prune}. The hypergraph $H$ is a set of distinct transversal $100$-subsets of $V_1\cup\dots\cup V_{100}$, so it is simple, $100$-partite and $100$-uniform on $100n$ vertices, and by \cref{prop:prune} every $99$-set lies in at most $200$ of its edges; this is (1).

For (2), let $c$ be a facet colouring of $H$. Each colour class is a family of pairwise facet-disjoint edges of $H\subseteq\Hrand$. Since $A$ holds, $\Hrand$ has no $m=T/5$ pairwise facet-disjoint edges, and deleting edges cannot create such a family, so every colour class has at most $m-1$ edges. If $c$ used at most $450$ colours, then
\[
|E(H)|\ \le\ 450(m-1)\ <\ 450\cdot\frac T5=90T,
\]
contradicting $|E(H)|>90T$. Hence $c$ uses at least $451$ colours.
\end{proof}

\phantomsection\label{pf:large}
\thmlarge*
\begin{proof}
Fix $d\ge15$ and put $a=\lceil5\log d\rceil$ and $x=a/d$. Since $g(t)=t-5\log t$ is increasing for $t\ge5$ and $g(15)>1.4$, we have $a\le5\log d+1<d$, so $x\in(0,1)$. Let $n$ be a multiple of $d$, $T=n^{d-1}$, and repeat the construction of \cref{sec:random} with $d$ classes of size $n$: $H_0$ is the complete $d$-partite $d$-uniform hypergraph, its facet hypergraph $F$ is $d$-uniform, $n$-regular and linear on $N=dT$ vertices, each edge is retained with probability $p=d/n$, and edges containing a $(d-1)$-set of codegree greater than $2d$ are deleted. Set $m=xN/d=xT=aT/d$, an integer since $d\mid n$ and $d-1\ge1$; $s=d\ge3$.

\emph{First moment.} By \cref{lem:count} (for $n$ large enough that $b=(d-1)(1-x)N\ge4$ and $\varepsilon_n=x/n+d^3/b<1$), exactly as in the proof of \cref{prop:firstmoment} with $c=pn=d$,
\[
\log\E[\#\text{$m$-matchings}]\ \le\ T\Bigl[x\log\frac dx+x-d\,f(x)\Bigr]+(m+u)(d+1)\varepsilon_n^{1/(d-1)} .
\]
Using $f(x)\ge x^2/2$ and $\log(d/x)=2\log d-\log a$, the bracket is at most
\[
x\Bigl[\log\frac dx+1-\frac{dx}2\Bigr]=x\Bigl[2\log d-\log a+1-\frac a2\Bigr]\ \le\ x\Bigl[1-\tfrac12\log d-\log a\Bigr]=:-\gamma_d ,
\]
where we used $a\ge5\log d$; and $\gamma_d>0$ because $\tfrac12\log d+\log a\ge\tfrac12\log15+\log14>1$. Since $\varepsilon_n\to0$ as $n\to\infty$ with $d$ fixed, there is $n_0(d)$ such that for $n\ge n_0(d)$ the error term is at most $\gamma_dT/2$, whence $\Pr(A^c)\le e^{-\gamma_dT/2}$, where $A$ is the event that $\Hrand$ has no $m$ pairwise facet-disjoint edges.

\emph{Pruning.} As in \cref{prop:prune}, with $X\sim\Bin(n-1,d/n)$, $\Pr(X\ge2d)\le2^{-2d}(1+d/n)^{n-1}\le(e/4)^d=:q_d$, so $\E Z\ge dT(1-dq_d)$, and $Z\le2dT$ deterministically. For $n\ge n_1(d)$ we have $2e^{-\gamma_dT/2}\le dq_d$ (as $T=n^{d-1}\to\infty$), and then
\[
\E[Z\ind_A]\ \ge\ dT(1-dq_d)-2dT\,e^{-\gamma_dT/2}\ \ge\ dT(1-2dq_d).
\]
So for $n\ge\max(n_0(d),n_1(d))$ some outcome satisfies $A$ and $|E(H)|\ge dT(1-2dq_d)$; let $H_d$ be such an $H$. It is simple, $d$-uniform, and every $(d-1)$-set lies in at most $2d$ of its edges.

\emph{Colouring.} As in the proof of \cref{thm:main}, every colour class of a facet colouring of $H_d$ has at most $m-1<xT$ edges, so
\[
\chif(H_d)\ \ge\ \frac{|E(H_d)|}{m-1}\ >\ \frac{dT(1-2dq_d)}{xT}=\bigl(1-2d(e/4)^d\bigr)\frac{d^2}{a} .
\]
Finally, with $r=2d$, $r+d-1=3d-1<3d$ and $a\le5\log d+1$, so $\chif(H_d)/(r+d-1)\ge(1-2d(e/4)^d)\,d/(3(5\log d+1))=(1-o(1))\,d/(15\log d)$, which tends to infinity. The limits were taken in the order ``fix $d$, then $n\to\infty$''; the thresholds $n_0(d),n_1(d)$ are not quantified here.
\end{proof}

\section{Report of the LLM referee (verbatim)}\label{app:referee}

\begin{tcolorbox}[colback=gray!8,colframe=gray!50,boxrule=0.4pt,arc=2pt]
\small The following is the unedited report of the adversarial referee agent (\texttt{claude-opus-5}, 2026-09-17) on the \emph{original} writeup. Its step and equation numbers refer to that writeup, not to this note (the writeup's (1), (3), (4), (6), (7), (9) are our \eqref{eq:count}, \eqref{eq:pm}, \cref{lem:int}, \eqref{eq:six}, \eqref{eq:Z} and the Stirling step of \cref{lem:count}). The scripts it mentions are in \texttt{verification\_astra/scripts/a\_generalization\_of\_vizings\_theorem/}. Treat it as machine output, not as an authority.
\end{tcolorbox}

\begingroup
\let\*\ast
\input{.build/referee.tex}
\endgroup

\begin{thebibliography}{99}
\small
\setlength{\itemsep}{2pt}
\bibitem{OPG} M.~Rosenfeld, A generalization of Vizing's Theorem?, Open Problem Garden, posted 11 April 2007, \url{http://www.openproblemgarden.org/op/a_generalization_of_vizings_theorem}.
\bibitem{Vizing64} V.~G. Vizing, On an estimate of the chromatic class of a $p$-graph, Diskret.\ Analiz 3 (1964) 25--30.
\bibitem{PS89} N.~Pippenger and J.~Spencer, Asymptotic behavior of the chromatic index for hypergraphs, J.~Combin.\ Theory Ser.~A 51 (1989) 24--42.
\bibitem{Kahn96} J.~Kahn, Asymptotically good list-colorings, J.~Combin.\ Theory Ser.~A 73 (1996) 1--59.
\bibitem{AK97} N.~Alon and J.~H. Kim, On the degree, size, and chromatic index of a uniform hypergraph, J.~Combin.\ Theory Ser.~A 77 (1997) 165--170.
\bibitem{Furedi86} Z.~F\"uredi, The chromatic index of simple hypergraphs, Graphs Combin.\ 2 (1986) 89--92.
\bibitem{Berge89} C.~Berge, On the chromatic index of a linear hypergraph and the Chv\'atal conjecture, Ann.\ New York Acad.\ Sci.\ 555 (1989) 40--44.
\bibitem{BFH24a} A.~Bretto, A.~Faisant and F.~Hennecart, The Berge--F\"uredi conjecture on the chromatic index of hypergraphs with large hyperedges, \href{https://arxiv.org/abs/2403.06850}{arXiv:2403.06850} (2024).
\bibitem{BFH24b} A.~Bretto, A.~Faisant and F.~Hennecart, About Berge--F\"uredi's conjecture on the chromatic index of hypergraphs, \href{https://arxiv.org/abs/2403.09518}{arXiv:2403.09518} (2024); C.~R.\ Math.\ Acad.\ Sci.\ Paris.
\bibitem{MA25} T.~Murff and X.~D. Arsiwalla, Upper bounds on the chromatic index of linear hypergraphs, \href{https://arxiv.org/abs/2510.07494}{arXiv:2510.07494} (2025).
\bibitem{KKKMO21} D.~Y. Kang, T.~Kelly, D.~K\"uhn, A.~Methuku and D.~Osthus, Solution to a problem of Erd\H{o}s on the chromatic index of hypergraphs with bounded codegree, \href{https://arxiv.org/abs/2110.06181}{arXiv:2110.06181} (2021).
\bibitem{Bregman73} L.~M. Br\'egman, Some properties of nonnegative matrices and their permanents, Soviet Math.\ Dokl.\ 14 (1973) 945--949.
\bibitem{Radhakrishnan97} J.~Radhakrishnan, An entropy proof of Bregman's theorem, J.~Combin.\ Theory Ser.~A 77 (1997) 161--164.
\bibitem{AF08} N.~Alon and S.~Friedland, The maximum number of perfect matchings in graphs with a given degree sequence, Electron.\ J.\ Combin.\ 15 (2008) N13.
\bibitem{CR11} J.~Cutler and A.~J. Radcliffe, An entropy proof of the Kahn--Lov\'asz theorem, Electron.\ J.\ Combin.\ 18 (2011) P10.
\end{thebibliography}

\end{document}
